% auto-generated by scripts/run_learner4.py -- do not edit by hand
Three strategies were compared at a fixed movement speed, so that the only thing varying is co-contraction: learn the full third-order plant without co-contracting (L1); co-contract and learn only the slow lag, treating the transients as noise (L2); and co-contract, then relax progressively and learn the whole model (L5, and L6 which warm-starts a full-order learner from the slow lag). On the slow-mode criterion $|\hat\tau_d-\tau_d|<0.08$\,s the co-contracting learners arrive in a median of 3 (L2), 3 (L5) and 3 (L6) trials against 6 for the soft learner, and they lose 0.0, 0.0 and 0.0 trials to failure against 4.0. Tested on a common soft, fast probe --- every condition asked to predict the same unbraced reach --- L5 reaches relative error 0.007 and L1 0.007, while the stiff slow-only learner plateaus at 0.216: co-contraction buys a correct slow model quickly, but only relaxation buys the transients. That the filter rather than the reduced parameterisation is responsible is shown by L4, a full-order learner held stiff, which ties with L2 on the slow criterion (3 trials), and by L3, a first-order learner held soft, which never reaches it and plateaus at +0.218\,s of bias.
